\chapter{Nightmares}

\section*{Definition}
Nightmares are vivid, disturbing dreams that cause awakening from sleep, with rapid return to alertness and clear recall of the frightening content. Nightmare disorder is diagnosed when they are frequent (at least weekly for 3 months) and cause clinically significant distress or impairment.

\section*{Pathophysiology}
Nightmares occur most commonly during REM sleep, which predominates in the second half of the night. They involve hyperactivation of the amygdala, medial prefrontal cortex, and anterior cingulate during REM, reflecting impaired fear memory extinction and exaggerated emotional reactivity in the dreaming brain.

\section*{Etiology}
\begin{itemize}
  \item \textbf{Psychiatric:} Post-traumatic stress disorder (most common pathological association), depression, anxiety disorders, borderline personality disorder
  \item \textbf{Trauma:} Combat, sexual assault, accidents, natural disasters, childhood abuse
  \item \textbf{Medication-induced:} Antidepressants (SSRIs, venlafaxine, bupropion -- especially during initial treatment), beta-blockers, levodopa, dopaminergic agonists
  \item \textbf{Substance withdrawal:} Alcohol, benzodiazepines, barbiturates, antidepressants
  \item \textbf{Sleep disorders:} Sleep apnea, restless legs syndrome, narcolepsy, REM sleep behavior disorder
\end{itemize}

\section*{Indications for Evaluation}
Seek care if:
\begin{itemize}
  \item Nightmares occur more than once per week and are associated with significant distress or sleep avoidance
  \item Nightmares began after a traumatic event or are associated with PTSD
  \item There is excessive daytime sleepiness, dream enactment behavior, or suspected sleep apnea
\end{itemize}

\section*{Related Symptoms}
\begin{itemize}
  \item PTSD, insomnia, sleep avoidance, fear of darkness, daytime fatigue, depression, anxiety, night terrors, sleep paralysis, REM sleep behavior disorder
\end{itemize}
\textbf{Note:} Nightmare disorder is distinct from night terrors (which occur in NREM sleep with no recall) and can persist for years without treatment.

\section*{Self-Care / Home Management}
\begin{itemize}
  \item Practice imagery rehearsal therapy: rewrite the nightmare with a different ending and rehearse the new version during the day
  \item Reduce alcohol and caffeine, especially before bed
  \item Establish a calming bedtime routine: warm bath, dim lights, relaxing activities
  \item Make the sleep environment feel safe (nightlight, door locked, stuffed animal for comfort)
\end{itemize}

\section*{Diagnostic Approach}
\begin{itemize}
  \item Clinical history assessing nightmare frequency, content, timing, and associated distress
  \item Screen for PTSD (PCL-5), depression (PHQ-9), and other sleep disorders
  \item Sleep diary for 2 weeks documenting bedtime, awakenings, nightmares, and daytime symptoms
  \item Polysomnography if comorbid sleep disorder suspected (sleep apnea, REM behavior disorder)
\end{itemize}

\section*{Treatment / Management Overview}
\begin{itemize}
  \item First-line psychotherapy: Imagery rehearsal therapy (IRT), with strong evidence for efficacy
  \item Exposure-based therapies for trauma-related nightmares (CBT for PTSD, EMDR)
  \item Pharmacotherapy: Prazosin (alpha-blocker, best evidence for trauma nightmares), clonidine, or trazodone
  \item Treat underlying PTSD, depression, or sleep disorders (CPAP for sleep apnea)
\end{itemize}

\clearpage
